\documentclass[11pt]{article}
\usepackage[margin=1in]{geometry}
\usepackage{amsmath,amssymb,amsthm,booktabs,graphicx,longtable}
\usepackage[colorlinks=true,urlcolor=blue,linkcolor=blue]{hyperref}
\newtheorem{proposition}{Proposition}
\newtheorem{example}{Example}
\newcommand{\E}{\mathbb E}
\newcommand{\V}{\operatorname{Var}}
\newcommand{\Cov}{\operatorname{Cov}}
\title{Causal inference with uncertain record linkage:\\a finite-population research compendium}
\author{Research draft}
\date{October 2026}
\begin{document}
\maketitle
\begin{abstract}
A linkage error rate does not determine causal bias. The fields moved by a false link,
which units are linked correctly, exclusion from the analysis, and the treatment
assignment mechanism all matter. We formalize these distinctions, prove finite-population
results, and implement two remedies: validation-sample correction and inference over
feasible reconstructions. The correction permits outcome-dependent linkage when an
independent probability sample validates the necessary fields. Graph ranges require
separate assumptions about candidate recall, overlap, and unmatched outcomes; ranges
of estimates alone are not causal confidence intervals. Exhaustive small cases check
the derivations, and simulations illustrate scope and failure. A read-only application
reconstructs the Rajasthan reservation--MNREGA linkage pipeline. The broad problem and
several remedies already have substantial precedents. This draft makes no claim that
the results establish a novel statistical method.
\end{abstract}

\section{What existing work establishes}
Causal inference with uncertain identities is an existing field. Propensity-score
subclassification methods already trade linkage confidence against downstream
estimation error \cite{wortman}. Bayesian approaches jointly infer identity and causal
quantities \cite{guha}, including covariates distributed across files \cite{guhax}.
Treatment-roster linkage has been analyzed with multiple imputation and an explicit
sensitivity analysis for informative linkage \cite{shan}. Experimental linking errors
can reduce statistical power \cite{tahamont}.

The closest broad comparator treats each of $X$, treatment, and outcome as the field
separated from the other two, derives naive-estimator bias, and develops corrections
with asymptotic inference \cite{slawski}. Its mismatch and auxiliary-variable assumptions
are conditional on covariates, not a blanket claim that all match quality is unrelated
to potential outcomes. Other methods allow analysis variables to inform linkage.
Consequently, ``existing work only keeps high-precision matches'' would misdescribe
the literature. These are statements about specified methods; no systematic survey of
applied practice was conducted here.

Optimization over feasible identity resolutions already bounds aggregate queries
\cite{turkcapar}. Optimization over acceptable treated--control pairings also exists
\cite{morucci}, although those pairings join different people, not records for the same
person. Validation-based correction belongs to the established two-phase estimation
literature \cite{twophase}. The contribution question here is narrower: whether retaining
the candidate graph, assignment restrictions, and joint field structure yields a useful
causal analysis beyond these methods. That question remains open. The literature matrix
records assumptions and exact source locations; absence from this matrix is not proof
of novelty.

\section{Target, files, and sources of randomness}
Fix $N$ target units $i=1,\ldots,N$, baseline covariates $X_i$, and potential outcomes
$Y_i(0),Y_i(1)$. Initially assume no interference, consistency, and perfect compliance.
Assignment $Z$ is uniform over all vectors with $m$ treated units, $0<m<N$, independent
of the fixed potential outcomes. Write $c=N-m$ and
\[
 Y_i=Z_iY_i(1)+(1-Z_i)Y_i(0),\qquad
 \tau_N=N^{-1}\sum_i\{Y_i(1)-Y_i(0)\}.
\]
The oracle difference in means is
\[
 \widehat\tau_O=\sum_iw_i(Z)Y_i,\qquad w_i(Z)=Z_i/m-(1-Z_i)/c.
\]
Its assignment expectation is $\tau_N$ and its variance is
$S_1^2/m+S_0^2/c-S_\tau^2/N$, where each $S^2$ uses divisor $N-1$.
The Neyman variance estimator omits the unidentifiable negative term. Its expectation
is conservative; this does not make a finite-sample normal interval exact.

File inclusion indicators $R_i^A,R_i^B$ determine which units have records. Noisy
identifiers $Q^A,Q^B$ generate a bipartite candidate graph $G$. A true partial mapping
$L^*$ links records for the same unit; it is injective only if each file has at most
one record per relevant unit and period. Multiple villages in a panchayat, boundary
changes, and repeated measurements must be resolved before imposing this assumption.
An estimated linkage $\widehat L$ and retention rule $S$ produce the analysis data.
The true link can be absent from $G$ even when a source has several candidates.

A reconstruction $D(L)$ moves all fields recorded together. Three useful layouts are
$(X,Z)\mid Y$, $(X,Y)\mid Z$, and $X\mid(Z,Y)$. Separately uncertain covariates and
outcomes require separate maps with explicit consistency restrictions. Reassigning
individual columns independently can create impossible records. We distinguish absent
counterparts, omitted true candidate edges, wrong accepted edges, ambiguous candidates,
donor reuse, and source exclusions. All may depend on $X$, assignment, realized outcomes,
or latent potential outcomes unless a result rules that dependence out.

Expectations must specify their randomness: assignment, population sampling, identifier
corruption, validation sampling, or an algorithm's random choices. Sections 3--5 use
fixed populations and state when averaging over validation or linkage as well. The
simulation regenerates populations, then assignments and errors. Conditioning on a
chosen linkage answers a different uncertainty question from repeating the entire
process.

With noncompliance, distinguish assignment $Z$, treatment receipt $D(Z)$, and outcomes
$Y(D,Z)$. The default target is the intention-to-treat effect. A ratio of outcome and
receipt contrasts targets a complier effect only with additional exclusion, monotonicity,
relevance, and design assumptions. Linkage uncertainty in both numerator and denominator
must be propagated jointly. A near-zero feasible first stage can make a ratio set
unbounded; dividing marginal endpoint bounds is generally invalid. Reservation status
and the sex of the elected officeholder are therefore not interchangeable treatments.

\section{How bias and variance enter}
\begin{proposition}[Bookkeeping identities, established]
For any two square-integrable estimators on the same probability space, define
$\Delta=\widehat\tau_L-\widehat\tau_O$. Then
\begin{align*}
 \E\widehat\tau_L-\tau_N&=(\E\widehat\tau_O-\tau_N)+\E\Delta,\\
 \V(\widehat\tau_L)&=\V(\widehat\tau_O)+\V(\Delta)+2\Cov(\widehat\tau_O,\Delta).
\end{align*}
\end{proposition}
\begin{proof}
Substitute $\widehat\tau_L=\widehat\tau_O+\Delta$, and use linearity of expectation
and the variance of a sum. No independence is needed.
\end{proof}
The covariance cannot be discarded. A procedure returning zero under a zero-effect
population can have less variance than the oracle. A smaller reported standard error
can coexist with greater bias. For a quadratic-loss decomposition,
$\operatorname{MSE}_L=\V(\widehat\tau_L)+(\E\widehat\tau_L-\tau_N)^2$.
A bootstrap of fixed linked rows neither discovers omitted true links nor automatically
accounts for the negative dependence created by donor exclusivity.

To separate mechanisms, let $a_O(S),r_O(S)$ be clean adjusted and unadjusted estimators
on a selected set, and $a_L(S),r_L(S)$ their reconstructed versions. Then
\[
 a_L(S)-a_O(N)=\underbrace{a_O(S)-a_O(N)}_{\text{selection}}
 +\underbrace{r_L(S)-r_O(S)}_{\text{reconstruction}}
 +\underbrace{[a_L(S)-r_L(S)]-[a_O(S)-r_O(S)]}_{\text{adjustment distortion}}.
\]
This identity requires all estimators to exist. Empty arms or failed fits are failures,
not zeros. A different order gives a different allocation of interactions; these are
not uniquely identified causal components of linkage error.

\begin{proposition}[Pre-assignment selection]
For a fixed retained subset $S$, conditional on its containing both assignment arms,
its unadjusted contrast has expectation $\tau_S=|S|^{-1}\sum_{i\in S}\tau_i$.
Writing $s_i=1\{i\in S\}$ and $\bar s=|S|/N$,
$\tau_S-\tau_N=\Cov_N(s_i,\tau_i)/\bar s$, where $\Cov_N$ uses divisor $N$.
\end{proposition}
\begin{proof}
Conditional on the treated count within $S$, all subsets of that size are equally
likely. Each arm mean is unbiased for its corresponding potential-outcome mean in
$S$. Averaging over nonempty counts proves the first assertion. Expand the covariance
for the second.
\end{proof}
Thus perfect precision does not remove target-population bias. If $S$ depends on realized
assignment or outcomes, this proof fails and even the retained-population contrast
need not be causally identified. Inclusion weighting needs positive, identified
inclusion probabilities and the appropriate selection ignorability; it cannot recover
units that have no chance of being observed without extrapolation.

\begin{proposition}[Exact finite-population permutation expectation]
Suppose outcome linkage is a fixed permutation $\pi$, independent of assignment.
Set $F_i=1\{\pi(i)=i\}$ and $q=N^{-1}\sum_iF_i$. Then
\begin{align*}
 \E_Z\sum_i w_i(Z)Y_{\pi(i)}
 &=\frac{\sum_i F_i\tau_i-\tau_N}{N-1}\\
 &=\frac{Nq-1}{N-1}\tau_N+
   \frac{N}{N-1}\Cov_N(F_i,\tau_i).
\end{align*}
\end{proposition}
\begin{proof}
For $i=j$, $\E\{w_iY_i\}=\tau_i/N$. For $i\ne j$,
$\E w_i=0$ and $\E(w_iZ_j)=-1/[N(N-1)]$, obtained from
$\Pr(Z_i=Z_j=1)=m(m-1)/[N(N-1)]$. Therefore
$\E(w_iY_j)=-\tau_j/[N(N-1)]$. Each outcome appears exactly once.
Summing diagonal and off-diagonal terms proves the first line; expanding the covariance
proves the second. This holds for unequal arm sizes as well.
\end{proof}
For random permutations independent of assignment, average the formula over their law;
$F_i$ is replaced by its correctness probability in the first line. Homogeneous effects
or equal correctness probabilities give a scalar attenuation factor. Without either
restriction, average precision alone is insufficient. A fixed derangement gives
$-\tau_N/(N-1)$, a small sign reversal rather than literal zero. This finite-population
calculation is a derivation supplied here, with no priority claim.

\begin{example}[Same precision, different causal information]
A within-arm permutation preserves the realized difference in means exactly, however
many identities are wrong. Swapping entire treatment and control outcome vectors in
an equally allocated experiment reverses its sign. Both can have zero correct links.
The first map depends on assignment, so applying the preceding independent-permutation
formula to it would be an error. Moving $Z$ and $Y$ jointly also preserves their
unadjusted contrast; moving only $X$ can still damage observational adjustment.
\end{example}

For outcome contamination let $e_z$ be the wrong-link fraction in observed arm $z$,
$\mu_z$ its correctly linked outcome mean, and $\nu_z$ the corresponding donor mean.
The linked contrast equals
$(1-e_1)\mu_1+e_1\nu_1-(1-e_0)\mu_0-e_0\nu_0$.
This is a mixture identity, not an identification theorem. Attenuation requires
restrictions on the donor means. When $\mu_1=1,\mu_0=0,e_1=e_0=.1$, choosing
$\nu_1=10,\nu_0=0$ gives 1.9; choosing $\nu_1=-20$ gives $-1.1$.
False links can amplify or reverse an effect.

For a different model, suppose treatment is randomized with probability $p$ and
misclassified nondifferentially: $a=\Pr(\widetilde Z=0\mid Z=1)$,
$b=\Pr(\widetilde Z=1\mid Z=0)$, with label error independent of potential outcomes
conditional on $Z$. Bayes' rule gives
\[
 \E[Y\mid\widetilde Z=1]-\E[Y\mid\widetilde Z=0]
 =\left\{\frac{p(1-a)}{p(1-a)+(1-p)b}
 -\frac{pa}{pa+(1-p)(1-b)}\right\}\tau.
\]
The proof substitutes $\E[Y\mid Z=z]=\E[Y(z)]$ into each posterior mixture.
A roster with fixed treated count induces dependencies not described by independent
Bernoulli misclassification. It needs the full assignment or linkage model.

\section{Validation as a remedy}
The following is a standard two-phase difference-estimation construction specialized
to a completely randomized experiment. It is useful because it does not require an
ignorable linkage mechanism.

\begin{proposition}[Validation correction and conservative variance expectation]
Let $h_i$ be any fully observed linked score, possibly constructed using all linked
records and assignments, with $H=\sum_i h_i$. Let an independent simple random sample
$A$ of $k\ge2$ anchor units validate both $Y_i$ and $Z_i$, after all $h_i$ are fixed.
Define $d_i=w_i(Z)Y_i-h_i$. Then
\[
 \widehat\tau_V=H+\frac Nk\sum_{i\in A}d_i,
 \qquad \E_A(\widehat\tau_V\mid Z,\text{files})=\widehat\tau_O.
\]
If both assignment arms have at least two units,
\[
 \widehat V_V=\widehat V_{O,A}+N^2(1-k/N)s_{d,A}^2/k
\]
has expectation $\V_{Z,A}(\widehat\tau_V)+S_\tau^2/N$, conditional on the finite
population and averaging over any linkage randomness, where
\[
 \widehat V_{O,A}=\frac1{\pi_2}\left\{
 \frac{\sum_{i<j\in A:Z_i=Z_j=1}(Y_i-Y_j)^2}{m^2(m-1)}+
 \frac{\sum_{i<j\in A:Z_i=Z_j=0}(Y_i-Y_j)^2}{c^2(c-1)}\right\},
 \quad \pi_2=\frac{k(k-1)}{N(N-1)}.
\]
\end{proposition}
\begin{proof}
The validation sample mean of $d$ is unbiased for its population mean. Its conditional
variance is $(1-k/N)S_d^2/k$, and $s_{d,A}^2$ is unbiased for $S_d^2$.
The conditional mean of the validation error is zero, so its covariance with the
oracle estimator is zero. The law of total variance gives the result once the oracle
variance is estimated. Each pair has inclusion probability $\pi_2$.
Using $\sum_{i<j}(Y_i-Y_j)^2=n\sum_i(Y_i-\bar Y)^2$ within each arm, the expectation
of $\widehat V_{O,A}$ over validation equals the full-data Neyman variance estimator.
Its assignment expectation exceeds the oracle variance by $S_\tau^2/N$.
\end{proof}
This result allows errors in treatment as well as outcomes, and arbitrary dependence
of linkage on either. It requires a complete anchor population, known assignment
counts, error-free validation of the necessary fields, and a probability validation
sample independent of the fixed scores. Reviewing only doubtful links violates this
sampling assumption. Scores can be zero for unlinked units, but validation must still
be capable of recovering their true outcomes. If records are intrinsically unavailable,
that requirement is not met. Targeted review needs its own inclusion probabilities;
the implemented formula cannot be reused unchanged. Normal intervals are assessed by
simulation, not guaranteed by the variance result.

A model-based alternative divides a linked contrast by
$\widehat\kappa=(N\widehat q-1)/(N-1)$. Even when the model justifies the population
factor, estimating and inverting it creates finite-sample bias. The delta variance is
\[
 V(\widehat T/\widehat\kappa)\simeq V(\widehat T)/\kappa^2+
 T^2V(\widehat\kappa)/\kappa^4-2T\Cov(\widehat T,\widehat\kappa)/\kappa^3.
\]
With fixed permutation and an independent validation sample, the covariance term is
zero under assignment--validation averaging. A random permutation or shared fitted
linkage model can invalidate that simplification. The simulation comparator uses a
working linked-data Neyman variance and is labeled approximate. Near-zero factors are
reported as failures. Fitting a mixture or joint Bayesian model is another established
route; the present code does not reproduce every estimator in the literature matrix.

\section{Feasible reconstructions and causal uncertainty}
For each source $i$ and donor $j$, let $x_{ij}\in\{0,1\}$ indicate a link and let
$u_i\in\{0,1\}$ indicate an outside option. The implemented feasible set imposes
\[
 \sum_jx_{ij}+u_i=1,\quad \sum_ix_{ij}\le1,\quad
 K_{\min}\le\sum_{ij}x_{ij}\le K_{\max},\quad
 \sum_{(i,j)\notin G}x_{ij}\le B_G.
\]
Only explicitly permitted sources may take the outside option. Accepted ``certain''
links can be imposed with a budget $B_C$ for violations. These are sensitivity
parameters, not probabilities estimated from a score margin. A missing candidate edge
and a genuinely absent outcome are different events. An off-graph budget permits
known donors outside the graph; an outside option requires outcome support for an
unobserved or unmatched outcome. Both may be needed.

\begin{proposition}[Sharp ranges of a linear reconstructed statistic]
For fixed weights $w_i$, donor outcomes $y_j$, and independently unconstrained missing
outcomes $v_i\in[a,b]$, the sharp lower and upper values of
$\sum_{ij}w_i y_jx_{ij}+\sum_iw_iv_iu_i$ over the stated feasible set are binary
linear programs. Outside coefficients are $\min(w_i a,w_i b)$ for the lower problem
and $\max(w_i a,w_i b)$ for the upper problem.
\end{proposition}
\begin{proof}
For any fixed matching, the outside terms minimize or maximize separately at the
appropriate endpoints. The remaining objective is linear in binary assignment
variables. Every feasible reconstruction corresponds to one feasible vector of these
variables, and each feasible vector constructs a permitted reconstruction with the
chosen endpoints. The finite optimum is attained, proving sharpness relative to these
assumptions. Extra joint restrictions on missing outcomes could narrow the range.
\end{proof}
Without budgets or other non-network constraints this reduces to an assignment/min-cost
flow problem. General budgets can require integer programming. ``Sharp'' here describes
the range of a statistic over specified reconstructions. It does not establish sharp
partial identification of $\tau_N$, because unobserved potential outcomes remain.

For $m$ treated-roster records that must match injectively into a complete population
of $N$ outcome records, all units outside the roster image are controls. Therefore
\[
 \widehat\tau_L=\frac{N}{m(N-m)}\sum_{j\in\operatorname{image}(L)}y_j
                  -\frac{\sum_jy_j}{N-m}.
\]
The denominators remain fixed. Incomplete roster coverage breaks this expression.
Ordinary least-squares coefficients with a fixed design matrix are also linear in
outcomes; moving treatment or covariates changes the design matrix, and generally
requires optimizing the entire estimator rather than holding its weights fixed.

\begin{example}[Outside options and global influence]
Two sources both have candidates with outcomes 0 and 10. Requiring both to match gives
mean 5. Allowing either to be unmatched with support $[0,10]$ gives mean range $[0,10]$.
Presence of candidates does not justify the first cardinality restriction.
Next, let source 1 have candidates $b_1=0,b_2=0$, and source 2 have
$b_1=0,b_3=100$, with both required to match. The mean range is $[0,50]$.
Learning that source 1 uses $b_1$ collapses the mean to 50; learning it uses $b_2$ leaves
the range unchanged. A zero local outcome range can thus carry review value because
of competition for donors. `review\_priority()` reports best and guaranteed reductions,
without assigning probabilities to possible reviewer answers.
\end{example}

\begin{proposition}[Coverage and graph failure are separate]
Suppose a full-data confidence set $C(D^*)$ covers a fixed causal target with probability
at least $1-\alpha$, and the random reconstruction set $\mathcal D$ contains $D^*$
with probability at least $1-\delta$. Then
$\bigcup_{D\in\mathcal D}C(D)$ covers with probability at least $1-\alpha-\delta$.
For valid full-data sharp-null $p$-values, $p_{\max}=\sup_{D\in\mathcal D}p(D)$
rejects at level $\alpha$ with probability at most $\alpha+\delta$.
\end{proposition}
\begin{proof}
Failure can occur only if the true reconstruction is excluded or the full-data
procedure fails at the truth. Apply the union bound. On the containment event,
$p_{\max}\ge p(D^*)$, proving the test statement.
\end{proof}
No independence between graph construction and the experiment is needed for this
unconditional argument. It does require that $C(D^*)$ or $p(D^*)$ be a valid oracle
procedure on the full-data probability space. If an analyst freezes an outcome-dependent
selected graph and permutes only the selected rows using a fictional assignment law,
that oracle premise is generally false. Either reproduce the selection/linkage process
under each assignment or derive the correct conditional law. Our exact test enumerates
the original complete assignment support, with all anchors retained, for every feasible
full linkage. It tests the sharp zero-effect null; it is not a test of the weak zero-ATE
null under heterogeneous effects.

\begin{proposition}[An explicit conservative causal confidence set]
Suppose every potential outcome lies in a known interval $[a,b]$, all anchors are
retained, and assignment is complete with counts $m,c$. Let $[L,U]$ bound the oracle
realized contrast over feasible reconstructions, and set
\[
 r=(b-a)\sqrt{\log(4/\alpha)/2}\{m^{-1/2}+c^{-1/2}\}.
\]
Then $[L-r,U+r]\cap[-(b-a),b-a]$ covers $\tau_N$ with probability at least
$1-\alpha-\delta$, where $\delta$ bounds graph/reconstruction failure.
\end{proposition}
\begin{proof}
Each observed arm mean is a simple-random-sample mean of its fixed potential-outcome
population. Hoeffding's bound for sampling without replacement \cite{hoeffding} gives, for each arm,
absolute deviation at most $(b-a)\sqrt{\log(4/\alpha)/(2n_z)}$ except on probability
$\alpha/2$. The triangle inequality and union bound control the difference. Apply the
previous proposition and intersect with the logical ATE support.
\end{proof}
These intervals can be wide; that is the cost of weak assumptions. The observed sample
minimum and maximum cannot substitute for known potential-outcome support. More
powerful design-specific intervals can be unioned, but their uniform validity must be
established separately. A credible interval from a linkage model answers a different
question and requires the specified posterior model.

\section{Observational extension and limits of identification}
With observational assignment, even perfect identity does not identify an ATE without
causal assumptions. Under consistency, conditional exchangeability
$\{Y(0),Y(1)\}\perp Z\mid X$, and positivity, an oracle estimator can use
$e(X)=\Pr(Z=1\mid X)$ and outcome regressions. If the linked covariates are
$\widetilde X$, exchangeability given $X$ does not imply exchangeability given
$\widetilde X$. A linkage model must preserve the conditioning information or quantify
its loss. Standard doubly robust estimation protects against misspecification of one
nuisance model on correctly measured data; it does not automatically protect against
incorrect identities shared by both nuisance fits.

\begin{proposition}[Correct measured-data models need not remove confounding]
Let $X$ be Bernoulli with probability $1/2$. Set
$\Pr(Z=1\mid X)=.2+.6X$, $Y(0)=X$, and $Y(1)=X+1$. Observe $W=X$ with probability $q$ and an
independent Bernoulli$(1/2)$ draw otherwise, with this mechanism independent of
$(X,Z)$ apart from the copying. Even using the true measured-covariate propensity
$e_W(w)=.5+.6q(w-.5)$, the expectation of its inverse-propensity estimator is
\[
 1+\frac{.6(1-q^2)}{1-.36q^2},
\]
whereas the causal effect and oracle $X$-adjusted expectation both equal one.
\end{proposition}
\begin{proof}
Bayes' rule gives $\Pr(X=1\mid W=w)=.5+q(w-.5)$, hence
$\V(X\mid W)=(1-q^2)/4$. Within either $W$ stratum,
$\Cov(X,Z\mid W)=.6\V(X\mid W)$ and
$e_W(1-e_W)=(1-.36q^2)/4$. The conditional difference in mean baseline outcomes
is the ratio of these last two quantities. Adding the treatment effect and averaging
over $W$ proves the expression. The same observed-data contrast is obtained with
correctly specified measured-data outcome regressions, so double robustness does not
repair the lost exchangeability.
\end{proof}
This replacement model describes an external donor distribution, not a fixed
one-to-one permutation. The separate script \texttt{scripts/observational.R}
checks the formula at five values of $q$ with known propensities. Here $q$ is a
copying probability; an incorrect donor can coincidentally have the same binary
covariate, so $q$ is not observed covariate agreement. This is an illustration of
residual confounding, not a novelty claim or a general observational remedy.

The bookkeeping identity still holds, but its oracle-bias term can be nonzero.
The validation difference principle applies to known-propensity oracle scores, with
the appropriate sampling variance. Estimated propensity scores and outcome models
need joint influence-function, cross-fitting, or resampling calculations that include
validation and linkage estimation. The finite-randomization variance formula above
cannot simply be copied. Likewise, a union over linkages inherits causal coverage only
from a valid observational oracle procedure. Linkage sensitivity does not bound
unmeasured confounding.

\begin{proposition}[Uninformative identities can prevent causal identification]
Consider two anchor units with exactly one assigned treatment. File A records their
identities and assignments. File B records the unordered pair of realized outcomes,
with no informative identifier; its candidate graph is complete. Even though every
unit appears once and assignment is randomized, the joint distribution of these
observations does not identify $\tau_2$.
\end{proposition}
\begin{proof}
In population P, both units have $Y_i(0)=0,Y_i(1)=1$, so $\tau_2=1$.
In population Q, both have $Y_i(0)=1,Y_i(1)=0$, so $\tau_2=-1$.
Under either population and either possible assignment, the unordered outcome file
is exactly $\{0,1\}$. File A's two possible assignment vectors each have probability
$1/2$ under both populations. Identifiers and the candidate graph are the same.
Thus the entire observable law is identical while the causal target differs.
\end{proof}
The unordered-file condition can equivalently be represented by sorting outcome rows
or randomly ordering them independently of identity. In that representation the
correct row mapping need not be independent of assignment, which is permitted here
but was excluded in the fixed-permutation theorem. This example proves more than two
compatible realized contrasts: it constructs identical observable probability laws.
Candidate containment and full overlap alone do not identify the effect. Validation,
informative identifiers, or restrictions on the linkage mechanism add information.
Randomization still supports conservative set inference with the known outcome support.

\section{Simulation and empirical illustration}
The executable protocol in \texttt{research/simulation-design.md} specifies populations,
corruption mechanisms, methods, seeds, and reporting before the pilot. The main
simulation crosses field layout and linkage mechanism, includes unequal treatment
allocation and heterogeneous effects, and separates controlled candidate graphs from
synthetic identifier corruption. The primary target stays the full finite population;
selected-target shifts are reported separately. Normal-approximation coverage is
measured rather than assumed. The graph confidence sets use the preceding bounded-outcome
result, and their widths should be judged alongside their protection.
\begin{table}[htbp]\centering\small
\caption{Finite-population ATE estimation with uncertain linkage. Each cell uses independent simulated populations of 80 units; validation methods audit 40. Bias is in outcome units. Coverage is for approximate 95\% Wald intervals; MCSE is Monte Carlo standard error.}
\begin{tabular}{llrrr}\toprule
Scenario & Method & Bias & Coverage & MCSE\\\midrule
independent y & naive & -0.304 & 0.795 & 0.009\\
independent y & attenuation & 0.003 & 0.959 & 0.004\\
independent y & audit difference & -0.004 & 0.945 & 0.005\\
covariate y & naive & -0.471 & 0.574 & 0.011\\
covariate y & attenuation & -0.237 & 0.893 & 0.007\\
covariate y & audit difference & -0.002 & 0.940 & 0.005\\
assignment y & naive & -0.005 & 0.954 & 0.005\\
assignment y & attenuation & 0.436 & 0.831 & 0.008\\
assignment y & audit difference & 0.005 & 0.948 & 0.005\\
independent t & naive & -0.300 & 0.809 & 0.009\\
independent t & attenuation & 0.013 & 0.954 & 0.005\\
independent t & audit difference & 0.015 & 0.942 & 0.005\\
selection & naive & -0.673 & 0.339 & 0.011\\
selection & attenuation & -0.526 & 0.745 & 0.010\\
selection & audit difference & -0.018 & 0.943 & 0.005\\
identifiers & naive & -0.122 & 0.921 & 0.006\\
identifiers & attenuation & -0.000 & 0.954 & 0.005\\
identifiers & audit difference & -0.003 & 0.942 & 0.005\\
\bottomrule\end{tabular}
\end{table}

\begin{table}[htbp]\centering\small
\caption{Point intervals and graph confidence sets, with 2000 replications per cell. Width is in outcome units. Null rejection uses sharp-zero effect runs; power uses positive-effect runs. Coverage refers to the finite-population ATE. Graph sets use known potential-outcome support [0,6].}
\begin{tabular}{llrrrr}\toprule
Scenario & Method & Coverage & Width & Null rejection & Power\\\midrule
identifiers & audit difference & 0.942 & 1.12 & 0.055 & 0.942\\
identifiers & graph & 1.000 & 12.00 & 0.000 & 0.000\\
identifiers & naive & 0.921 & 1.01 & 0.057 & 0.925\\
independent y & audit difference & 0.945 & 1.28 & 0.053 & 0.860\\
independent y & graph & 1.000 & 7.00 & 0.000 & 0.000\\
independent y & naive & 0.795 & 1.04 & 0.057 & 0.751\\
selection & audit difference & 0.943 & 2.29 & 0.057 & 0.385\\
selection & graph & 1.000 & 11.26 & 0.000 & 0.000\\
selection & naive & 0.339 & 1.10 & 0.047 & 0.222\\
\bottomrule\end{tabular}\end{table}

Under independent outcome linkage, the naive estimator has bias -0.304; the validation difference estimator has bias -0.004 (MCSE 0.007). Under selection, their biases are -0.673 and -0.018. Validation restores the fixed-population target under its sampling assumptions; the normal approximation still needs separate assessment.

The graph causal sets are protective but weak in these designs. For independent outcome linkage their mean width is 7.00 and power is 0.000; validation intervals have width 1.28 and power 0.860. The identifier design permits every source an outside option, so its graph sets can be logically uninformative. The separate estimator ranges still measure how much one-to-one coupling constrains reconstruction. This experiment does not establish a competitive general-purpose graph interval.


The Rajasthan application is implemented independently in
\texttt{application/run.R}. It hashes source inputs, reproduces the current accepted
links and regression specification, preserves all candidates before nearest-neighbor
selection, and compares exact and fuzzy linkage on both method-specific and common
cohorts. Its source audit separates reservation assignment from officeholder sex and
records the missing historical allocation information. The application is a linkage
and estimator sensitivity exercise; it does not identify actual linkage bias without
independently validated identities. General legal provisions for lotteries do not
establish the year-specific randomization support needed for an exact test.
\begin{table}[htbp]\centering\small
\caption{Rajasthan completed projects, 2011--2014: reproduced 2005 reservation coefficient, controlling for 2010 reservation. Standard errors condition on linked data. Cohort changes differ from linkage reconstruction uncertainty.}
\begin{tabular}{llrrr}\toprule
Link rule & Cohort & GPs & Estimate & OLS SE\\\midrule
fuzzy & method & 4355 & -0.026 & 1.880\\
fuzzy & common & 1861 & -2.992 & 2.902\\
exact & method & 2036 & -2.994 & 2.740\\
exact & common & 1861 & -2.992 & 2.902\\
\bottomrule\end{tabular}
\end{table}
On the fixed 4,355-GP fuzzy cohort, candidate assignment extrema for the 2005 coefficient are $[-29.72,31.89]$, compared with $[-48.60,47.27]$ when donor reuse is allowed. These are conditional estimator ranges, not causal confidence intervals.


\section{Result status and next empirical requirements}
The bias/variance decomposition, selection weighting, mixture calculations,
two-phase correction, graph optimization, and union-of-confidence-sets argument use
established tools. The exact finite-population permutation and validation-variance
formulas are proved and tested here, but their priority is unresolved. The globally
coupled review example is a useful diagnostic, not by itself a paper contribution.

The implemented remedies require different information. A probability validation
sample buys design-based correction without outcome-independent linking. A credible
candidate set and outcome support buy sensitivity intervals without a calibrated
linkage probability model. An informative joint model can provide narrower estimates
at the cost of modeling assumptions. The application needs independent identity checks,
historical assignment rosters and lottery rules, and a dated geographic crosswalk
before it can support stronger causal or linkage-bias claims. Comparing all contemporary
model-based estimators, extending valid inference to observational nuisance estimation,
and developing efficient sharp-null optimization for large graphs remain research
extensions rather than completed claims.

\begingroup\small
\begin{thebibliography}{20}
\bibitem{wortman} Wortman and Reiter. Simultaneous record linkage and causal inference
with propensity score subclassification. 2018.
\url{https://arxiv.org/abs/1709.03631}.
\bibitem{guha} Guha, Reiter, and Mercatanti. Bayesian causal inference with bipartite
record linkage. 2022 (preprint 2020). \url{https://arxiv.org/abs/2002.09119}.
\bibitem{guhax} Guha and Reiter. Regression-assisted Bayesian record linkage for causal
inference in observational studies with covariates spread over two files. 2024.
\url{https://doi.org/10.1016/j.jspi.2023.07.004}.
\bibitem{shan} Shan, Thomas, and Gutman. A multiple imputation procedure for record
linkage and causal inference to estimate the effects of home-delivered meals. 2021.
\url{https://pmc.ncbi.nlm.nih.gov/articles/PMC9222523/}.
\bibitem{tahamont} Tahamont et al. Dude, where's my treatment effect? Errors in
administrative data linking and the destruction of statistical power in randomized
experiments. 2021. \url{https://doi.org/10.1007/s10940-020-09461-x}.
\bibitem{slawski} Slawski. Causal secondary analysis of linked data in the presence of
mismatch error. 2025 preprint. \url{https://arxiv.org/abs/2512.14492}.
\bibitem{turkcapar} Turkcapar and Krishnan. Quantifying uncertainty in aggregate queries
over integrated datasets. 2023. \url{https://arxiv.org/abs/2309.05178}.
\bibitem{morucci} Morucci, Noor-E-Alam, and Rudin. A robust approach to quantifying
uncertainty in matching problems of causal inference. 2022.
\url{https://doi.org/10.1287/ijds.2022.0020}.
\bibitem{twophase} Amorim et al. Two-phase sampling designs for data validation in
settings with covariate measurement error and continuous outcome. 2021.
\url{https://pmc.ncbi.nlm.nih.gov/articles/PMC8715909/}.
\bibitem{hoeffding} Hoeffding. Probability inequalities for sums of bounded random
variables. Journal of the American Statistical Association, 58:13--30, 1963.
\url{https://doi.org/10.1080/01621459.1963.10500830}.
\end{thebibliography}
\endgroup
\end{document}
